\begin{abstract}

The Prescription Drug User Fee Act of 1992 (PDUFA) and the Generic Drug User Fee Amendments of 2012 (GDUFA) authorized the FDA to collect industry fees in exchange for binding review-time targets. Standard incentive logic predicts that user-fee acceleration should disproportionately benefit drugs with the largest expected revenue streams, including many controlled substances. I test this hypothesis using the Drugs@FDA database linked to DEA scheduling records (1939--2025, $N = 25{,}908$ original approvals) and estimate interrupted time series, Poisson rate models with exposure, and stacked difference-in-differences specifications. Once group-specific trends are allowed, the estimated effect on the controlled-substance share is indistinguishable from zero for both policies (PDUFA NDA-only IRR $= 1.19$, 95\% CI $[0.59, 2.39]$; GDUFA ANDA-only IRR $= 1.38$, 95\% CI $[0.86, 2.21]$). A separable late-2000s wave of branded abuse-deterrent opioid NDAs explains apparent rate movements in less flexible specifications.

\end{abstract}
